\documentclass[12pt,a4paper]{article}
\usepackage{fontspec}
\setmainfont{Liberation Serif}
\renewcommand{\contentsname}{}
\usepackage[margin=2.2cm]{geometry}
\usepackage{amsmath}
\usepackage{booktabs}
\usepackage{array}
\usepackage{tikz}
\usepackage{enumitem}
\usepackage{xcolor}
\usepackage{titlesec}
\setlength{\parindent}{0pt}
\sloppy
\setlength{\parskip}{6pt}
\titleformat{\section}{\large\bfseries}{}{0pt}{}
\titleformat{\subsection}{\normalsize\bfseries}{}{0pt}{}

\newcommand{\gantt}[2]{%
\begin{center}
\begin{tikzpicture}
\pgfmathsetmacro{\gu}{14/#1}
\node[below,font=\scriptsize] at (0,0) {0};
\foreach \n/\s/\e in {#2}{
  \draw[fill=gray!12] (\s*\gu,0) rectangle (\e*\gu,0.75);
  \node[font=\small] at ({(\s+\e)/2*\gu},0.375) {\n};
  \node[below,font=\scriptsize] at (\e*\gu,0) {\e};
}
\end{tikzpicture}
\end{center}}

\newcommand{\q}[1]{\par\noindent\textbf{Асуулт.} #1\par}
\newcommand{\ans}{\par\noindent\textbf{Хариулт.} }

\begin{document}

\begin{center}
{\Large\bfseries Лабораторийн ажил №5}\\[4pt]
{\large Процесс ба санах ой: 5-р бүлэг. CPU-ийн хуваарилалт (CPU Scheduling)}\\[4pt]
Practice Exercises 5.1--5.10\\[10pt]
\begin{tabular}{rl}
Оюутан: & Ч. Мөнхбаяр \\
Код: & B242270045 \\
Хичээл: & Үйлдлийн систем \\
Ном: & Silberschatz, Galvin, Gagne. \textit{Operating System Concepts}, 10th ed.
\end{tabular}
\end{center}

\vspace{6pt}
\hrule
\vspace{6pt}

Энэхүү тайланд сурах бичгийн 5-р бүлгийн Practice Exercises хэсгийн 5.1--5.10 даалгаврын асуултыг монгол хэлээр орчуулж, тус бүрийн хариултыг бичив. Тооцоотой бодлогуудад Гантын диаграм, эргэлтийн хугацаа (turnaround time) болон хүлээлтийн хугацааг (waiting time) хүснэгтээр харуулав.

Тайланд дараах тодорхойлолтуудыг ашиглана:
\begin{itemize}[nosep]
\item Эргэлтийн хугацаа = Дууссан хугацаа $-$ Ирсэн хугацаа
\item Хүлээлтийн хугацаа = Эргэлтийн хугацаа $-$ CPU burst хугацаа
\end{itemize}

%-------------------------------------------------
\section{Даалгавар 5.1}
\q{CPU хуваарилалтын алгоритм нь хуваарилагдсан процессуудын гүйцэтгэх дарааллыг тодорхойлдог. Нэг процессор дээр $n$ процессыг хуваарилах бол хичнээн өөр хуваарь боломжтой вэ? $n$-ээр илэрхийлсэн томьёо гарга.}
\ans Процесс бүрийг таслалгүй (нэг удаа бүтнээр нь) гүйцэтгэнэ гэж үзвэл энэ нь $n$ процессын бүх сэлгэмэлийн тоо болно. Эхний байранд $n$ процессын аль нэг, хоёр дахь байранд үлдсэн $n-1$-ийн аль нэг гэх мэтээр сонгогдох тул
\[
N = n \cdot (n-1) \cdot (n-2) \cdots 2 \cdot 1 = n!
\]
Жишээлбэл, 3 процесс байвал $3! = 6$ өөр хуваарь боломжтой.

%-------------------------------------------------
\section{Даалгавар 5.2}
\q{Таслах (preemptive) болон таслахгүй (nonpreemptive) хуваарилалтын ялгааг тайлбарла.}
\ans
\textbf{Таслахгүй хуваарилалт} үед процесс CPU-г авсан бол өөрөө суллах хүртлээ, өөрөөр хэлбэл дуусах эсвэл I/O хүлээх (waiting төлөвт шилжих) хүртлээ ашиглана. Үйлдлийн систем түүнийг хүчээр зогсоохгүй. Хэрэгжүүлэхэд хялбар, контекст сэлгэлт цөөн боловч урт процесс богино процессуудыг удаан хүлээлгэж болно.

\textbf{Таслах хуваарилалт} үед үйлдлийн систем ажиллаж буй процессоос CPU-г хүчээр авч өөр процесст өгч болно. Жишээ нь, илүү өндөр эрэмбэтэй процесс ирэх эсвэл time quantum дуусах үед таймерын тасалдлаар процессыг ready дараалал руу буцаана. Энэ нь хариу өгөх хугацааг сайжруулж, интерактив болон бодит цагийн системд тохиромжтой. Харин контекст сэлгэлтийн нэмэлт зардал гарч, хуваалцсан өгөгдөл дээр race condition үүсэхээс сэргийлэх синхрончлол шаардлагатай болно.

%-------------------------------------------------
\section{Даалгавар 5.3}
\q{Дараах процессууд заасан хугацаанд ирнэ. Таслахгүй хуваарилалт ашиглаж, шийдвэр гаргах тухайн мөчид мэдэгдэж буй мэдээлэл дээр үндэслэнэ.}
\begin{center}
\begin{tabular}{ccc}
\toprule
Процесс & Ирэх хугацаа & Burst хугацаа \\
\midrule
$P_1$ & 0.0 & 8 \\
$P_2$ & 0.4 & 4 \\
$P_3$ & 1.0 & 1 \\
\bottomrule
\end{tabular}
\end{center}
\begin{enumerate}[label=\alph*.,nosep]
\item FCFS алгоритмын үеийн дундаж эргэлтийн хугацаа хэд вэ?
\item SJF алгоритмын үеийн дундаж эргэлтийн хугацаа хэд вэ?
\item Эхний 1 нэгж хугацаанд CPU-г сул байлгаад дараа нь SJF хэрэглэвэл дундаж эргэлтийн хугацаа хэд болох вэ? (future-knowledge scheduling)
\end{enumerate}

\ans

\textbf{a. FCFS.} 0 мөчид зөвхөн $P_1$ ирсэн тул түүнийг ажиллуулж, дараа нь ирсэн дарааллаар $P_2$, $P_3$ ажиллана.
\gantt{14}{$P_1$/0/8, $P_2$/8/12, $P_3$/12/13}
\begin{center}
\begin{tabular}{cccc}
\toprule
Процесс & Ирэх & Дуусах & Эргэлтийн хугацаа \\
\midrule
$P_1$ & 0.0 & 8 & 8.0 \\
$P_2$ & 0.4 & 12 & 11.6 \\
$P_3$ & 1.0 & 13 & 12.0 \\
\bottomrule
\end{tabular}
\end{center}
Дундаж эргэлтийн хугацаа $= (8 + 11.6 + 12)/3 = 31.6/3 \approx \mathbf{10.53}$.

\textbf{b. SJF.} 0 мөчид зөвхөн $P_1$ байгаа тул түүнийг сонгоно. Таслахгүй тул $P_1$ 8 хүртэл ажиллана. 8 мөчид $P_2$ (4) болон $P_3$ (1) хүлээж байгаа тул богино нь болох $P_3$, дараа нь $P_2$.
\gantt{14}{$P_1$/0/8, $P_3$/8/9, $P_2$/9/13}
\begin{center}
\begin{tabular}{cccc}
\toprule
Процесс & Ирэх & Дуусах & Эргэлтийн хугацаа \\
\midrule
$P_1$ & 0.0 & 8 & 8.0 \\
$P_2$ & 0.4 & 13 & 12.6 \\
$P_3$ & 1.0 & 9 & 8.0 \\
\bottomrule
\end{tabular}
\end{center}
Дундаж эргэлтийн хугацаа $= (8 + 12.6 + 8)/3 = 28.6/3 \approx \mathbf{9.53}$.

\textbf{c. Эхний 1 нэгжид сул зогсоод SJF.} 1 мөчид бүх процесс ирсэн байх тул хамгийн богино $P_3$, дараа нь $P_2$, эцэст нь $P_1$ ажиллана.
\gantt{14}{сул/0/1, $P_3$/1/2, $P_2$/2/6, $P_1$/6/14}
\begin{center}
\begin{tabular}{cccc}
\toprule
Процесс & Ирэх & Дуусах & Эргэлтийн хугацаа \\
\midrule
$P_1$ & 0.0 & 14 & 14.0 \\
$P_2$ & 0.4 & 6 & 5.6 \\
$P_3$ & 1.0 & 2 & 1.0 \\
\bottomrule
\end{tabular}
\end{center}
Дундаж эргэлтийн хугацаа $= (14 + 5.6 + 1)/3 = 20.6/3 \approx \mathbf{6.87}$.

Эндээс харахад CPU-г 1 нэгж сул байлгасан ч ирээдүйн мэдээллийг ашигласнаар дундаж эргэлтийн хугацаа SJF-ээс ч бага гарч байна.

%-------------------------------------------------
\section{Даалгавар 5.4}
\q{Дараах процессуудын CPU burst хугацааг миллисекундээр өгөв. Процессууд $P_1, P_2, P_3, P_4, P_5$ дарааллаар бүгд 0 мөчид ирсэн.}
\begin{center}
\begin{tabular}{ccc}
\toprule
Процесс & Burst хугацаа & Эрэмбэ \\
\midrule
$P_1$ & 2 & 2 \\
$P_2$ & 1 & 1 \\
$P_3$ & 8 & 4 \\
$P_4$ & 4 & 2 \\
$P_5$ & 5 & 3 \\
\bottomrule
\end{tabular}
\end{center}
\begin{enumerate}[label=\alph*.,nosep]
\item FCFS, SJF, таслахгүй эрэмбийн (их тоо нь өндөр эрэмбэ), RR (quantum = 2) алгоритмуудын Гантын 4 диаграм зур.
\item Алгоритм тус бүрийн үеийн процесс бүрийн эргэлтийн хугацаа хэд вэ?
\item Алгоритм тус бүрийн үеийн процесс бүрийн хүлээлтийн хугацаа хэд вэ?
\item Аль алгоритм хамгийн бага дундаж хүлээлтийн хугацаатай вэ?
\end{enumerate}

\ans

\textbf{a. Гантын диаграмууд.}

FCFS:
\gantt{20}{$P_1$/0/2, $P_2$/2/3, $P_3$/3/11, $P_4$/11/15, $P_5$/15/20}

SJF (тэнцүү burst байхгүй):
\gantt{20}{$P_2$/0/1, $P_1$/1/3, $P_4$/3/7, $P_5$/7/12, $P_3$/12/20}

Таслахгүй эрэмбэ ($P_1$ ба $P_4$-ийн эрэмбэ тэнцүү тул ирсэн дарааллаар буюу FCFS-ээр):
\gantt{20}{$P_3$/0/8, $P_5$/8/13, $P_1$/13/15, $P_4$/15/19, $P_2$/19/20}

RR (quantum = 2):
\gantt{20}{$P_1$/0/2, $P_2$/2/3, $P_3$/3/5, $P_4$/5/7, $P_5$/7/9, $P_3$/9/11, $P_4$/11/13, $P_5$/13/15, $P_3$/15/17, $P_5$/17/18, $P_3$/18/20}

\textbf{b. Эргэлтийн хугацаа} (бүгд 0-д ирсэн тул дууссан хугацаатай тэнцүү):
\begin{center}
\begin{tabular}{ccccc}
\toprule
Процесс & FCFS & SJF & Эрэмбэ & RR \\
\midrule
$P_1$ & 2 & 3 & 15 & 2 \\
$P_2$ & 3 & 1 & 20 & 3 \\
$P_3$ & 11 & 20 & 8 & 20 \\
$P_4$ & 15 & 7 & 19 & 13 \\
$P_5$ & 20 & 12 & 13 & 18 \\
\midrule
Дундаж & 10.2 & 8.6 & 15.0 & 11.2 \\
\bottomrule
\end{tabular}
\end{center}

\textbf{c. Хүлээлтийн хугацаа} (эргэлтийн хугацаа $-$ burst):
\begin{center}
\begin{tabular}{ccccc}
\toprule
Процесс & FCFS & SJF & Эрэмбэ & RR \\
\midrule
$P_1$ & 0 & 1 & 13 & 0 \\
$P_2$ & 2 & 0 & 19 & 2 \\
$P_3$ & 3 & 12 & 0 & 12 \\
$P_4$ & 11 & 3 & 15 & 9 \\
$P_5$ & 15 & 7 & 8 & 13 \\
\midrule
Дундаж & 6.2 & 4.6 & 11.0 & 7.2 \\
\bottomrule
\end{tabular}
\end{center}

\textbf{d.} Хамгийн бага дундаж хүлээлтийн хугацаатай нь \textbf{SJF} (4.6 мс). Энэ нь SJF нь дундаж хүлээлтийн хугацааг хамгийн бага болгодог оновчтой алгоритм гэдэгтэй нийцэж байна.

%-------------------------------------------------
\section{Даалгавар 5.5}
\q{Дараах процессуудыг таслах, round-robin хуваарилалтаар хуваарилна. Процесс бүр тоон эрэмбэтэй бөгөөд их тоо нь өндөр эрэмбийг илтгэнэ. Эрэмбэ нь 0 бөгөөд CPU нөөц хэрэглэдэггүй $P_{idle}$ процесс өөр ажиллах процесс байхгүй үед ажиллана. Time quantum нь 10 нэгж. Өндөр эрэмбэтэй процесс таславал тасалдсан процесс дарааллын төгсгөлд очно.}
\begin{center}
\begin{tabular}{cccc}
\toprule
Процесс & Эрэмбэ & Burst & Ирэх \\
\midrule
$P_1$ & 40 & 20 & 0 \\
$P_2$ & 30 & 25 & 25 \\
$P_3$ & 30 & 25 & 30 \\
$P_4$ & 35 & 15 & 60 \\
$P_5$ & 5 & 10 & 100 \\
$P_6$ & 10 & 10 & 105 \\
\bottomrule
\end{tabular}
\end{center}
\begin{enumerate}[label=\alph*.,nosep]
\item Процессуудын хуваарилалтын дарааллыг Гантын диаграмаар үзүүл.
\item Процесс бүрийн эргэлтийн хугацаа хэд вэ?
\item Процесс бүрийн хүлээлтийн хугацаа хэд вэ?
\item CPU-ийн ашиглалтын хувь хэд вэ?
\end{enumerate}

\ans

\textbf{a.} Хуваарилалтын явц:
\begin{itemize}[nosep]
\item 0--20: зөвхөн $P_1$ байгаа тул quantum дуусахад ч дахин сонгогдож 20-д дуусна.
\item 20--25: ажиллах процесс байхгүй, $P_{idle}$.
\item 25--35: $P_2$ ирж ажиллана. 30-д ижил эрэмбийн $P_3$ ирэх боловч таслахгүй, $P_2$ quantum-аа дуусгана.
\item 35--45: $P_3$; 45--55: $P_2$ (үлдэгдэл 5).
\item 55--60: $P_3$. 60-д эрэмбэ 35 бүхий $P_4$ ирж $P_3$-г таслана. $P_3$ дарааллын төгсгөлд очно (үлдэгдэл 10).
\item 60--75: $P_4$ бүтнээр ажиллаж дуусна.
\item 75--80: $P_2$ дуусна; 80--90: $P_3$ дуусна.
\item 90--100: $P_{idle}$.
\item 100--105: $P_5$. 105-д эрэмбэ 10 бүхий $P_6$ ирж таслана.
\item 105--115: $P_6$ дуусна; 115--120: $P_5$ дуусна.
\end{itemize}
\gantt{120}{$P_1$/0/20, idle/20/25, $P_2$/25/35, $P_3$/35/45, $P_2$/45/55, $P_3$/55/60, $P_4$/60/75, $P_2$/75/80, $P_3$/80/90, idle/90/100, $P_5$/100/105, $P_6$/105/115, $P_5$/115/120}

\textbf{b, c.} Эргэлтийн болон хүлээлтийн хугацаа:
\begin{center}
\begin{tabular}{cccccc}
\toprule
Процесс & Ирэх & Burst & Дуусах & Эргэлт & Хүлээлт \\
\midrule
$P_1$ & 0 & 20 & 20 & 20 & 0 \\
$P_2$ & 25 & 25 & 80 & 55 & 30 \\
$P_3$ & 30 & 25 & 90 & 60 & 35 \\
$P_4$ & 60 & 15 & 75 & 15 & 0 \\
$P_5$ & 100 & 10 & 120 & 20 & 10 \\
$P_6$ & 105 & 10 & 115 & 10 & 0 \\
\bottomrule
\end{tabular}
\end{center}

\textbf{d.} Нийт 120 нэгж хугацаанаас CPU сул байсан хугацаа $(25-20) + (100-90) = 15$ нэгж. Иймд
\[
\text{CPU ашиглалт} = \frac{120 - 15}{120} = \frac{105}{120} = \mathbf{87.5\%}.
\]

%-------------------------------------------------
\section{Даалгавар 5.6}
\q{Олон түвшний дарааллын системийн түвшин бүрт өөр өөр хэмжээтэй time quantum байх нь ямар давуу талтай вэ?}
\ans Процессууд шинж чанараараа ялгаатай тул түвшин бүрт тохирсон quantum өгөх нь үр ашгийг нэмэгдүүлнэ. Байнга I/O хийдэг интерактив процессууд богино CPU burst-тай тул тэднийг жижиг quantum-тай дээд түвшинд байлгавал хурдан хариу өгч, ready дараалалд удаан хүлээхгүй. Харин CPU-д түшиглэсэн урт процессууд том quantum-тай доод түвшинд байрласнаар контекст сэлгэлт цөөрч, CPU-ийн цаг илүү ашигтай зарцуулагдана. Ингэснээр нэг системд интерактив процессын хариу өгөх хурд ба урт процессын гүйцэтгэлийн үр ашгийг зэрэг хангах боломжтой.

%-------------------------------------------------
\section{Даалгавар 5.7}
\q{Олон CPU хуваарилалтын алгоритм параметртэй байдаг. Жишээ нь, RR алгоритм time slice параметр шаарддаг бол олон түвшний буцах холбоост дараалал (MLFQ) дарааллын тоо, дараалал бүрийн алгоритм, процессыг дараалал хооронд шилжүүлэх шалгуур зэрэг параметр шаарддаг. Иймд эдгээр нь үнэндээ алгоритмын олонлог юм. Нэг олонлог нөгөөгөө агуулж болно (жишээ нь, FCFS бол хязгааргүй quantum-тай RR). Дараах хосуудын хооронд ямар хамаарал байх вэ?}
\begin{enumerate}[label=\alph*.,nosep]
\item Эрэмбэ ба SJF
\item MLFQ ба FCFS
\item Эрэмбэ ба FCFS
\item RR ба SJF
\end{enumerate}
\ans
\begin{enumerate}[label=\alph*.]
\item \textbf{Эрэмбэ ба SJF.} SJF нь эрэмбийн алгоритмын тусгай тохиолдол. Дараагийн CPU burst хамгийн богино процесс хамгийн өндөр эрэмбэтэй байхаар (эрэмбэ = таамагласан burst-ийн урвуу) тохируулбал эрэмбийн хуваарилалт SJF болно. Иймд SJF $\subset$ Эрэмбэ.
\item \textbf{MLFQ ба FCFS.} FCFS нь MLFQ-ийн тусгай тохиолдол. Ганцхан дараалалтай, тэр дараалал нь FCFS (эсвэл хязгааргүй quantum-тай RR) бөгөөд процесс дараалал хооронд шилждэггүй MLFQ нь яг FCFS болно. Иймд FCFS $\subset$ MLFQ.
\item \textbf{Эрэмбэ ба FCFS.} FCFS нь эрэмбийн алгоритмын тусгай тохиолдол. Хамгийн эрт ирсэн процесс хамгийн өндөр эрэмбэтэй (эрэмбэ нь ирсэн хугацаанаас хамаарна) байвал эрэмбийн алгоритм FCFS шиг ажиллана. Иймд FCFS $\subset$ Эрэмбэ.
\item \textbf{RR ба SJF.} Хооронд нь агуулах хамаарал байхгүй. RR нь burst-ийн уртыг огт харгалздаггүй, процесс бүрт ижил quantum өгдөг бол SJF нь burst-ийн уртаар сонголт хийдэг. Аль ч quantum-ийн утгад RR нь SJF болохгүй, мөн эсрэгээрээ.
\end{enumerate}

%-------------------------------------------------
\section{Даалгавар 5.8}
\q{CPU хуваарилалтын алгоритм нь сүүлийн үед хамгийн бага процессорын хугацаа ашигласан процессуудыг илүүд үздэг гэж үзье. Энэ алгоритм яагаад I/O-д түшиглэсэн програмуудыг илүүд үзэх боловч CPU-д түшиглэсэн програмуудыг байнга өлсгөхгүй вэ?}
\ans I/O-д түшиглэсэн програмууд CPU-г богино хугацаанд ашиглаад удалгүй I/O хүлээх төлөвт шилждэг тул сүүлийн үеийн CPU хэрэглээ нь бага байдаг. Иймд тэд I/O-оос буцаж ready төлөвт ороход өндөр эрэмбэ авч, түрүүлж сонгогдоно.

Гэвч CPU-д түшиглэсэн програм байнга өлсөхгүй. Учир нь I/O-д түшиглэсэн програмууд CPU-г цөөхөн ашиглаад байнга сулладаг тул CPU-д чөлөөт хугацаа гарч, тэр үед CPU-д түшиглэсэн програм ажиллана. Мөн алгоритм зөвхөн \emph{сүүлийн үеийн} хэрэглээг харгалздаг. CPU-д түшиглэсэн програм ажиллахгүй хүлээх тусам түүний сүүлийн үеийн CPU хэрэглээ буурч, эрэмбэ нь аажмаар өсөх тул эцэст нь заавал CPU авна. Энэ нь хөгшрөлт (aging)-тэй төстэй үйлчилгээ үзүүлнэ.

%-------------------------------------------------
\section{Даалгавар 5.9}
\q{PCS болон SCS хуваарилалтыг ялга.}
\ans
\textbf{PCS (Process-Contention Scope)} буюу процессын өрсөлдөөний хүрээ нь нэг процессын доторх хэрэглэгчийн түвшний thread-үүд тухайн процесст хуваарилагдсан LWP (kernel thread) дээр ажиллахын төлөө хоорондоо өрсөлддөг хуваарилалт юм. Үүнийг thread-ийн сан (library) хийдэг бөгөөд many-to-one, many-to-many загварт хэрэглэгдэнэ. Ихэвчлэн программист тогтоосон эрэмбээр сонгоно.

\textbf{SCS (System-Contention Scope)} буюу системийн өрсөлдөөний хүрээ нь систем дэх бүх kernel thread-үүд физик CPU-ийн төлөө өрсөлддөг хуваарилалт бөгөөд үүнийг үйлдлийн системийн цөм (kernel) гүйцэтгэнэ. One-to-one загвартай Windows, Linux зэрэг систем зөвхөн SCS ашигладаг.

Товчоор, PCS нь процесс доторх thread-үүдийн хоорондох, SCS нь бүх системийн thread-үүдийн хоорондох өрсөлдөөн юм. Pthreads-д үүнийг \texttt{PTHREAD\_SCOPE\_PROCESS} болон \texttt{PTHREAD\_SCOPE\_SYSTEM} утгаар тохируулна.

%-------------------------------------------------
\section{Даалгавар 5.10}
\q{Уламжлалт UNIX хуваарилагч нь эрэмбийн тоо ба эрэмбийн хооронд урвуу хамаарал тогтоодог: тоо их байх тусам эрэмбэ бага. Хуваарилагч секунд тутамд эрэмбийг дараах томьёогоор дахин тооцоолно:
\[
\text{Эрэмбэ} = \frac{\text{сүүлийн CPU хэрэглээ}}{2} + \text{суурь},
\]
энд суурь $= 60$. $P_1$-ийн сүүлийн CPU хэрэглээ 40, $P_2$-ынх 18, $P_3$-ынх 10 бол дахин тооцоолсны дараах шинэ эрэмбүүд ямар байх вэ? Үүнд үндэслэн уламжлалт UNIX хуваарилагч CPU-д түшиглэсэн процессын харьцангуй эрэмбийг өсгөх үү, бууруулах уу?}
\ans
\begin{align*}
P_1 &: \; 40/2 + 60 = \mathbf{80}\\
P_2 &: \; 18/2 + 60 = \mathbf{69}\\
P_3 &: \; 10/2 + 60 = \mathbf{65}
\end{align*}
Тоо бага байх тусам эрэмбэ өндөр тул эрэмбийн дараалал $P_3 > P_2 > P_1$ болно. CPU-г хамгийн их ашигласан $P_1$ хамгийн том тоо буюу хамгийн бага эрэмбэ авсан. Иймд уламжлалт UNIX хуваарилагч CPU-д түшиглэсэн процессын харьцангуй эрэмбийг \textbf{бууруулдаг}. Ингэснээр CPU-г бага ашигладаг интерактив болон I/O-д түшиглэсэн процессууд илүүд үзэгдэнэ.

%-------------------------------------------------
\section{Дүгнэлт}
Энэхүү лабораторийн ажлаар CPU хуваарилалтын үндсэн ойлголт болох таслах ба таслахгүй хуваарилалт, FCFS, SJF, эрэмбийн болон RR алгоритмуудыг Гантын диаграмаар загварчилж, эргэлтийн болон хүлээлтийн хугацааг тооцоолов. 5.4-р даалгаварт SJF нь хамгийн бага дундаж хүлээлтийн хугацаа (4.6 мс) өгсөн нь онолтой нийцэж байна. Мөн олон түвшний дараалал, PCS ба SCS, уламжлалт UNIX хуваарилагчийн эрэмбийн тооцоолол зэргээр бодит системүүд I/O-д түшиглэсэн процессыг хэрхэн илүүд үзэж, CPU-д түшиглэсэн процессыг өлсгөхгүй байлгадаг талаар судлав.

\end{document}
